% ============================================================================
% DRAFT v2 -- Section 1: Introduction
% Genre: benchmark / evaluation study (not a new-detector paper).
% Genre precedents cited: Bouthillier (variance in ML benchmarks), IM-IAD
% (AD evaluation protocol), Yu et al. 2020 (NAS vs random), Cao et al.
% (streaming AD benchmark), Yu et al. 2026 (event-level evaluation).
%
% Numbers trace to proxy_diag/samplesize_analysis.json, phaseB_truth.json and
% docs/03_实验记录.md. Checked with adbench/numsnap.py after each edit.
% ============================================================================

\section{Introduction}
\label{sec:introduction}

Visual inspection on a production line has to work without defect samples.
Collecting and annotating the failures of a manufacturing process is expensive,
and the failure modes that matter are by definition the ones not yet seen, so the
established approach is to model normality and flag deviations. Reverse
distillation~\cite{deng2022reverse} is a representative and widely deployed member
of this family: a frozen teacher extracts features at several scales, a student is
trained to reproduce them on normal images, and at test time the discrepancy
between the two is taken as the anomaly signal.

Two engineering decisions have to be made before such a system can be deployed.
The first is the \emph{architecture}: what network the student should be. The
second is the \emph{scoring policy}: how the per-pixel discrepancy maps are
combined into the single number that decides whether a part is rejected. The
second decision is usually made implicitly --- the published pipeline averages the
discrepancy map over all pixels, a choice that is rarely stated and, to our
knowledge, never justified. The literature has overwhelmingly invested in the
first, applying neural architecture search (NAS) to anomaly
detection~\cite{liu2023simplenet,wang2021stpm} on the argument that the optimal
architecture is dataset-dependent and should be discovered rather than chosen.

\runinhead{The margins the field argues about, and the resolution of its
benchmark.} Progress on the standard benchmark is reported in tenths of a
percentage point. PatchCore reached $99.1\%$ image-level AUROC on MVTec-AD in
2022~\cite{roth2022patchcore}; the strongest methods reported since reach
$99.6\%$~\cite{liu2023simplenet,guo2025dinomaly}. The contested margin is
therefore $0.5$ points, on a benchmark whose test sets hold roughly a hundred
images per category. Before asking which design decisions produce such margins,
it is worth asking whether the standard evaluation design can resolve them at all.
We measure that first, because the answer determines how every subsequent
comparison in this literature should be read.

This paper is accordingly a \emph{benchmark-and-evaluation study}, not a new
detector. Work of this kind measures the evaluation practice of a field rather
than adding a method, and it has clear precedents in this area. Yu et
al.~\cite{yu2020evaluating} compared architectures returned by several NAS
algorithms against architectures drawn uniformly at random from the same search
space, and found the search phase uninformative. IM-IAD~\cite{xie2024imiad}
consolidated competing evaluation protocols for industrial anomaly detection onto
one footing. Bouthillier et al.~\cite{bouthillier2021accounting} showed that
unacknowledged run-to-run variance invalidates comparisons in machine-learning
benchmarks and derived how many runs such comparisons require; the same
re-examination has since been carried out for
streaming~\cite{cao2024revisiting} and industrial time-series~\cite{yu2026eventlevel}
anomaly detection. We continue that line into a deployed unsupervised
image-anomaly-detection pipeline: we instrument it, we re-run its search, and we
train the complete comparison it never ran.

\runinhead{What we instrument.}
Our own earlier pipeline, KG-NAS, is a three-stage system of this family: a search
ranks 42 configurations (14 backbones $\times$ 3 feature dimensions) on normal
data alone using a knowledge-distillation signal, and the winning configuration is
then trained as a reverse-distillation detector. We ask a question this literature
has largely assumed rather than measured: \emph{how much accuracy does the
architecture choice actually buy, is the search that selects it doing its job, and
what can the benchmark that judges both actually resolve?} We answer by
measurement: we instrument the search, we train all 14 backbones of its search
space on all 15 MVTec-AD categories under one protocol --- 210 detectors --- and
we sweep the scoring policy on frozen models.

\runinhead{The benchmark resolves $3$ points; the field argues about $0.5$.}
Two-way decomposition of the 210-detector matrix leaves a residual standard
deviation of $3.30$ points, and the run-to-run standard deviation of a single
detector at a fixed (backbone, category) is $1.96$ points over three seeds.
Comparing two backbones across the fifteen categories is therefore a paired test
whose minimum detectable difference is \textbf{$2.98$ points} --- larger than the
entire $3.98$-point spread of the fourteen backbones, and six times the margin
that separates the 2022 and 2026 state of the art. To resolve $0.5$ points on a
single category would take $240$ independent training runs per backbone. This is
the paper's measurement about the field as a whole rather than about this
pipeline.

\runinhead{Two thirds of the search space cannot reach the deployed model.}
The searched feature dimension enters the search through the student projector, so
the search signal depends on it, but it does not reach the detector: in the
deployment path the student encoder builds its bottleneck with a hard-coded width
while the decoder is constructed with the requested one, and the final training
entry point never passes the argument at all. The three values of the searched
dimension therefore produce the identical deployed model for a given backbone. The
42-configuration space contains 14 distinct models; the other 28 cells are
deployment-equivalent duplicates. The recorded searches behave as this predicts:
across four independent runs the selected backbone is the same every time, while
the selected dimension differs (Section~\ref{subsec:inert}).

\runinhead{The search signal does not distinguish architectures, and collapses
for some of them.} Calling the search's own evaluation method ten times per
configuration gives two results. Across the twelve well-behaved backbones the
median score spans only $5.7\%$ --- the same order as the within-configuration
noise of a single backbone --- so the signal carries essentially no
architecture-level information. And two backbones fail in a specific way:
\texttt{efficientnet\_b1} returns exactly $0.0$ in $73\%$ of evaluations and
\texttt{mobilenetv3\_small\_100} in $13\%$, because the student diverges, every
validation batch trips a hard threshold, and the score is then silently computed
as an average over zero samples. Since the search minimises, $0.0$ is the best
score attainable anywhere in the space
(Section~\ref{subsec:proxy_stability}); over the project's own record of 834
search trials the collapse occurs in $8.0\%$ of them.

\runinhead{Architecture is detectable, negligible, and does not transfer.}
Training 210 detectors under an identical protocol and decomposing the resulting
AUROC matrix attributes $0.54\%$ of the variance to the backbone and $95.32\%$ to
the category, with the backbone effect statistically significant ($p = 0.039$) and
practically irrelevant. The manuscript's selected architecture ranks fifth of
fourteen and is statistically indistinguishable from twelve of the thirteen
alternatives. More consequentially, the ranking does not transfer: the mean
correlation between two categories' orderings of the 14 backbones is
$\rho = 0.215$, and in $31\%$ of category pairs the categories disagree. Yet the
protocol searches on a single category and applies the winner to all fifteen
(Section~\ref{subsec:capacity_null}).

\runinhead{The consequence is measurable, and it is worse than chance.}
Running the search itself with independent seeds, the two backbones that collapse
to zero are selected every time, and the architecture the search returns is worse
than picking a backbone uniformly at random. A design axis that cannot reach the
model, ranked by a signal that does not discriminate, on a quantity that does not
transfer, is selected --- reliably --- by a failure mode
(Section~\ref{subsec:search_vs_random}).

\runinhead{The leverage is in the decision nobody documents.} On frozen models,
changing only the reduction that turns discrepancy maps into an image score moves
mean AUROC from $82.36\%$ to $90.28\%$ --- $+7.92$ points, reproduced across three
training seeds --- and choosing the policy per category reaches $93.72\%$. The
policy can further be chosen without any anomaly label, by perturbing held-out
normal images with synthetic defects and selecting the policy that best separates
perturbed from clean: across three training seeds this recovers $7.09$ of the
$11.65$ points available from an oracle, on the same trained models
(Sections~\ref{subsec:policy_results} and~\ref{subsec:pseudo_results}).

\subsection*{Contributions}

\begin{enumerate}
  \item \textbf{The resolvability of the standard benchmark design.} We anchor the
  noise floor of this task to the margin the field actually contests, and show
  that a single-category, single-run comparison resolves $2.98$ points against a
  contested margin of $0.5$ --- so most published backbone comparisons on this
  benchmark are not measurements. This is a measurement about the evaluation
  practice of the field, and it is our primary contribution.
  \item \textbf{A diagnosis of the architecture axis.} Two thirds of a standard
  search space is deployment-inert; its search signal does not discriminate
  between architectures and degenerates to a score of exactly zero for a subset of
  them; the axis contributes $0.54\%$ of the variance in detection accuracy
  against $95.32\%$ for the category.
  \item \textbf{A controlled comparison of a full search space.} 14 backbones
  $\times$ 15 categories under one protocol (210 detectors), which makes the
  negative result measurable rather than anecdotal.
  \item \textbf{Evidence that a single-category search cannot work.} The backbone
  ranking does not transfer across categories ($\bar\rho = 0.215$; $31\%$ of
  category pairs reverse), so a search run on one category estimates a quantity
  that does not exist independently of the category.
  \item \textbf{A label-free method for the axis that does matter.} Selecting the
  scoring policy by separation of synthetic pseudo-anomalies recovers $61\%$ of the
  oracle gap --- $7.09$ of $11.65$ points, averaged over three training seeds ---
  without any anomaly label and without retraining.
\end{enumerate}

\subsection*{Scope and honesty}

We do not claim state of the art: the pipeline reaches $82.47\%$ mean AUROC on
MVTec-AD, against published values of $95.6\%$--$99.6\%$ for the established
methods, and Section~\ref{sec:positioning} reports that gap rather than obscuring
it. The contribution is not a better number but a quantitative account of which
design choices produce the number, of how much of the number the benchmark can
measure at all, and of the finding that the axis the field searches is not the
axis that matters. Section~\ref{sec:related_work} positions the work,
Section~\ref{sec:method} describes the pipeline under study,
Section~\ref{sec:diagnosis} reports the diagnosis and the resolvability analysis,
Section~\ref{sec:leverage} the scoring-policy analysis,
Section~\ref{sec:positioning} states where the pipeline stands, and
Section~\ref{sec:conclusion} concludes.
